This section provides an overview of the smart plant monitoring system, including its primary features, functions, and user interactions. The system is designed to monitor plant health using an integrated camera and environmental sensors, and to provide actionable insights through a user interface. From the perspective of end users, the system enables real-time monitoring and feedback. For maintainers and developers, it provides a structured backend for data processing and analysis. The key components, system behavior, and user interactions are described below.

\subsection{Features \& Functions}

The smart plant monitoring system is designed to capture plant images and environmental data, process this information, and provide users with plant health insights.

The system consists of the following core components:
\begin{itemize}
    \item \textbf{Camera Module:} Captures images of the plant at regular intervals.
    \item \textbf{Processing Unit:} Processes captured data and determines plant health status using backend logic.
    \item \textbf{Backend Server:} Stores data and performs analysis.
    \item \textbf{User Interface (Dashboard):} Displays plant health status, historical data, and alerts.
\end{itemize}

The system performs the following functions:
\begin{itemize}
    \item Captures plant images automatically
    \item Processes data to determine plant health status
    \item Stores data for historical tracking
    \item Displays results and alerts to the user
\end{itemize}

The system does \textbf{not} perform the following:
\begin{itemize}
    \item It does not physically interact with or modify the plant (no watering or actuation)
    \item It does not provide medical or agricultural guarantees
    \item It does not operate outdoors in extreme environmental conditions
\end{itemize}

External elements include the Internet (for data transfer) and optional cloud storage services.

\subsection{External Inputs \& Outputs}

The system receives input from both users and hardware components, and produces outputs for user interpretation.

\begin{center}
\begin{tabular}{|l|l|l|}
\hline
\textbf{Name} & \textbf{Description} & \textbf{Use} \\
\hline
Camera Image & Image of plant captured periodically & Used for plant health analysis \\
User Input & User settings (e.g., thresholds) & Configures system behavior \\
\hline
Health Status & Computed plant condition (e.g., healthy/unhealthy) & Displayed to user \\
Alerts & Notifications for abnormal conditions & Warns user of issues \\
Historical Data & Stored past readings and images & Used for tracking trends \\
\hline
\end{tabular}
\end{center}

\subsection{Product Interfaces}

The system provides a user-facing dashboard interface and internal system interfaces.

\textbf{User Interface:}
\begin{itemize}
    \item Displays plant health status (e.g., healthy, warning, critical)
    \item Provides historical graphs of plant conditions
    \item Displays alerts and notifications
\end{itemize}

\textbf{Hardware Interface:}
\begin{itemize}
    \item Camera connected to processing unit
    \item Sensors connected via microcontroller (e.g., Arduino or Raspberry Pi)
\end{itemize}

\textbf{Software Interface:}
\begin{itemize}
    \item Backend server processes incoming data
    \item Communication via APIs between hardware and backend
    \item Optional cloud database for storage
\end{itemize}